\chapter{Lernziele}
\begin{todolist}
    \item Ich kann die Prüfziffern eines \gls{EAN}-Codes berechnen bei gegebener Formel.
    \item Ich kann bestimmen, ob ein gegebener \gls{EAN}-Code korrekt ist.
    \item Ich kann bestimmen, welche Fehler (Vertippen einer Ziffer, Vertauschen von zwei Ziffern, Ziffer vergessen) durch eine Prüfziffer erkannt werden.
    \item Ich kann den Abstand zweier Bitfolgen gleicher Länge bestimmen.
    \item Ich kann den Mindest-Abstand einer Codierung bestimmen.
    \item Ich kann bestimmen, wie viele Fehler eine Codierung erkennt.
    \item Ich kann bestimmen, wie viele Fehler eine Codierung korrigieren kann.
    \item Ich kann erklären, welchen Abstand \textit{k}-fehlererkennende Codierungen haben müssen.
    \item Ich kann erklären, welchen Abstand \textit{k}-fehlerkorrigierende Codierungen haben müssen.
    \item Ich kann Codierungen angeben, die einen bestimmten Abstand haben, bzw. eine bestimmte Anzahl Fehler erkennen oder korrigieren können.
    \item Ich kann die \enquote{Kartentrick-Codierung} einer Bitfolge angeben.
    \item Ich kann erklären, wie die optimale Kartentrickcodierung (= optimale Länge und Breite für das Rechteck) für eine Bitfolge mit einer gegebenen Länge aussieht.
    \item Ich kann überprüfen, ob eine gegebene Kartentrick-Codierung korrekt ist.
    \item Ich kann bei einer Kartentrick-Codierung mit einem Fehler den Fehler finden und korrigieren.
    \item Ich kann in einem vorgegebenen Hyperwürfel (\textit{n}~$\leq 4$) die Zonen eines 1-fehlerkorrigierenden Codes einzeichnen.
    \item Ich kann an einem Hyperwürfel erklären, wann ein 1-fehlerkorrigierender Code perfekt ist.
    \item Ich kann den \acrshort{XOR}-Operator auf mehrere Bitfolgen gleicher Länge anwenden.
    \item Ich kann \gls{RAID}-0-, \gls{RAID}-1-, \gls{RAID}-5- und \gls{RAID}-6-Konfigurationen schematisch illustrieren und deren Nutzen erklären, indem ich die Begriffe \textit{striping}, \textit{mirroring} und \textit{Paritätsbits} verwende.
    \item Ich kann für eine vorgegebene Anzahl Festplatten sowie Anforderungen an Redundanz und Schreibgeschwindigkeit bestimmen, welche der vier \gls{RAID}-Konfigurationen am besten geeignet ist.
    \item Ich kann ein Paritätsbit mit dem \acrshort{XOR}-Operator berechnen.
    \item Ich kann erklären, weshalb \gls{RAID} kein Backup ist, und Szenarien nennen, in denen nur ein Backup die Daten rettet.
    \item Ich kann die 3-2-1-Regel und ihre Erweiterung 3-2-1-1-0 erklären und eine Backup-Strategie damit beurteilen.
\end{todolist}
